\section{引言：什么是 Harness？}

在所有工程领域中，harness（线束/挽具）都指代同一种东西：\textbf{连接、保护和编排各组件的基础设施层}——它本身并不执行核心工作。汽车中的线束（wiring harness）在引擎、传感器和仪表盘之间路由信号；测试线束（test harness）为代码提供可重复执行和可观测的脚手架；安全挽具（safety harness）则在你坠落时保护你。

Inngest 的这篇博客提出了一个核心论点：\textbf{Agent 运行时同样需要一个 harness}。LLM 是引擎，Tool 是外设，Memory 是存储——但连接它们的是什么？当 LLM 在第五次迭代超时时，谁来捕获这个失败？当两条消息冲突时，谁来防止竞态？当一个 webhook 事件到达时，谁来路由到正确的 handler？

\begin{importantbox}{核心论点：Harness vs. Framework}
当前每个 Agent 框架都在从零构建自己的重试逻辑、状态持久化、任务队列和事件路由。然而，\textbf{持久化的事件驱动基础设施}（durable, event-driven infrastructure）已经解决了这些问题。将每次 LLM 调用或 Tool 调用建模为一个 \textbf{step}——一个可独立重试的工作单元——就是 harness 思维的核心。
\end{importantbox}

文章的核心观点可以用一句话概括：\textbf{基础设施本身就是 harness}。框架关注的是 AI 逻辑的抽象，而 harness 关注的是让 AI 逻辑在生产环境中\textbf{可靠运行}的基础设施。

\subsection{Harness 与 Framework 的区别}

为了更清晰地理解二者的区别，下表从多个维度进行对比：

\begin{table}[H]
\centering
\caption{Harness 与 Framework 的对比}
\begin{tabular}{lll}
\toprule
\textbf{维度} & \textbf{Framework（框架）} & \textbf{Harness（线束）} \\
\midrule
关注点 & AI 逻辑的抽象与组合 & 可靠运行的基础设施 \\
重试机制 & 自行实现 & 内置 step 级重试 \\
状态持久化 & 自行实现 & 自动持久化每个 step 结果 \\
并发控制 & 自行实现 & 声明式 singleton 配置 \\
可观测性 & 需要额外集成 & step 级 trace 内置 \\
事件路由 & 通常不涉及 & 核心能力 \\
故障恢复 & 从头重跑或自定义逻辑 & 从上次成功的 step 恢复 \\
\bottomrule
\end{tabular}
\end{table}

\begin{knowledgebox}{类比理解：为什么叫 ``Harness''？}
考虑汽车的线束（wiring harness）：它不是引擎，不是轮胎，但没有它，引擎的信号无法传递到仪表盘，传感器的数据无法到达 ECU。线束解决的是\textbf{连接和保护}的问题。同理，Agent harness 解决的不是``如何调用 LLM''的问题，而是``LLM 调用失败后怎么办''``多个请求同时到达怎么办''这类\textbf{基础设施层面}的问题。
\end{knowledgebox}

\subsection{本章小结}

本章引入了 harness 的概念，指出当前 Agent 框架的一个共性问题：它们都在重复造轮子——重试、持久化、并发、事件路由——这些本质上是基础设施问题，不是 AI 问题。Inngest 的核心主张是用已有的 durable execution 基础设施来充当 Agent 的 harness。


